\section*{Zusammenfassung}
In dieser Arbeit wird ein Framework zur Robotersimulation entwickelt, das an der Fachhochschule Dortmund in der Lehre eingesetzt wird.  
Das Framework ermöglicht das Teach-In von Gelenkkonfigurationen sowie die Animation und Steuerung von Industrierobotern mit serieller Kinematik,
die mithilfe der Vorwärts- und Rückwärtstransformation berechnet werden.  
Beliebige Robotermodelle können mithilfe des Unified Robot Description Format (URDF) in das Framework integriert werden, wobei auch Roboter mit 7 oder mehr Drehgelenken in dem Framework Simuliert werden können,
da bei der Rückwärtstransformation das Newton-Raphson-Verfahren verwendet wird.
Als mathematisches Modell dient die Denavit-Hartenberg-Konvention, deren Koordinatensysteme visuell dargestellt werden.  
Darüber hinaus unterstützt das Framework die Simulation und Animation mobiler Roboter.  
Interaktive Steuerungen, wie beispielsweise Schieberegler zur Rotation von Würfeln, veranschaulichen kinematische Bewegungen direkt im Frontend.  
Ziel des Frameworks ist eine anschauliche und praxisnahe Vermittlung kinematischer Konzepte in der Robotik.  
Ein durchgeführter Stresstest bestätigt, dass das Framework unter typischen Systemvoraussetzungen flüssig und stabil läuft.


